\documentclass[11pt]{article}
\usepackage[utf8]{inputenc}
\usepackage[T1]{fontenc}
\usepackage[a4paper,margin=1in]{geometry}
\usepackage{lmodern}
\usepackage{microtype}
\usepackage{amsmath,amssymb,mathtools}
\usepackage{mathrsfs}
\usepackage{tikz-cd}
\usepackage{enumitem}
\usepackage{longtable,booktabs,array}
\usepackage[hidelinks,hypertexnames=false]{hyperref}
\setlength{\parindent}{0pt}
\setlength{\parskip}{0.55em}
\emergencystretch=3em
\hbadness=10000
\newcommand{\K}{K}
\newcommand{\Gr}{\operatorname{Gr}}
\newcommand{\Fil}{\operatorname{Fil}}
\newcommand{\ch}{\operatorname{ch}}
\newcommand{\Todd}{\operatorname{Todd}}
\newcommand{\cl}{\operatorname{cl}}
\newcommand{\Parf}{\operatorname{Parf}}
\newcommand{\Tor}{\operatorname{Tor}}
\newcommand{\Spec}{\operatorname{Spec}}
\newcommand{\RRR}{\mathrm{RRR}}
\providecommand{\codim}{\operatorname{codim}}
\providecommand{\Inv}{\operatorname{Inv}}
\providecommand{\Pic}{\operatorname{Pic}}
\providecommand{\Z}{\mathbb Z}
\providecommand{\Q}{\mathbb Q}
\providecommand{\Ahat}{\widehat A}
\providecommand{\Ahp}{\widehat A^{+}}
\providecommand{\That}{\widehat T}
\providecommand{\rad}{\operatorname{rad}}
\providecommand{\Gl}{\operatorname{Gl}}
\providecommand{\GL}{\operatorname{GL}}
\providecommand{\ch}{\operatorname{ch}}
\providecommand{\Todd}{\operatorname{Todd}}
\providecommand{\Gr}{\operatorname{Gr}}

% Batch 066 macro additions
\providecommand{\Atil}{\widetilde A}
\providecommand{\Ctil}{\widetilde C}
\providecommand{\K}{\mathrm K}
\providecommand{\cC}{\mathcal C}
\providecommand{\cF}{\mathscr F}
\providecommand{\eps}{\varepsilon}
\providecommand{\ellc}{\ell}
\providecommand{\Tor}{\operatorname{Tor}}
\providecommand{\Coker}{\operatorname{Coker}}
\providecommand{\Ker}{\operatorname{Ker}}
\providecommand{\Supp}{\operatorname{Supp}}
\providecommand{\codim}{\operatorname{codim}}
\providecommand{\Gr}{\operatorname{G}}
\providecommand{\rank}{\operatorname{rank}}
\providecommand{\rang}{\operatorname{rang}}
\providecommand{\cO}{\mathcal O}

\begin{document}

\providecommand{\D}{\mathrm D}
\providecommand{\Kcat}{\mathrm K}
\providecommand{\Qcoh}{\operatorname{Qcoh}}
\providecommand{\Mod}{\operatorname{Mod}}
\providecommand{\Coh}{\operatorname{Coh}}
\providecommand{\Lib}{\operatorname{Lib}}
\providecommand{\Loclib}{\operatorname{Loclib}}
\providecommand{\Hom}{\operatorname{Hom}}
\providecommand{\Ext}{\operatorname{Ext}}
\providecommand{\Ob}{\operatorname{ob}}
\providecommand{\Spec}{\operatorname{Spec}}
\providecommand{\RHom}{\mathrm R\!\operatorname{Hom}}
\providecommand{\ilim}{\varinjlim}
\providecommand{\plim}{\varprojlim}
\providecommand{\qcoh}{\mathrm{qcoh}}
\providecommand{\Id}{\operatorname{Id}}
\begin{center}\Large SGA 6 -- Exposé II, pages 191--205\end{center}


\section*{Exposé II. Existence of global resolutions, continued}

\textbf{Proposition 2.4.8.} Let $S$ be a compact complex pseudo-analytic Stein space (2.4.7). Suppose that, for every open $U$ of $S$ and every reduced closed subspace $X$ of $U$, the normalization of $X$ is locally connected. This condition is verified, for example, if $S$ is locally defined by real analytic equations and inequalities in some $\mathbb R^n$. Then
\[
A=\Gamma(S,\mathcal O_S)
\]
is a noetherian ring.

\emph{Proof.} It is enough to show that, if $M$ is an $A$-module of finite presentation, the set of submodules of $M$ (equivalently, the set of finite type submodules of $M$) satisfies the maximal condition. If $N$ is a submodule of finite type of $M$, then $M/N$ is of finite presentation by I 2.6 and I 2.9; hence $N$ is of finite presentation by the equivalence (2.4.3 c)). Using once more the equivalence (2.4.3 c)) between coherent $\mathcal O_S$-modules and $A$-modules of finite presentation, and using compactness of $S$, we are reduced to the following lemma.

\textbf{Lemma 2.4.8.1.} Let $S$ be a complex pseudo-analytic space satisfying the local connectedness condition of (2.4.8). Let $F$ be a coherent $\mathcal O_S$-module, let $(E_n)$ be an increasing sequence of coherent sub-$\mathcal O_S$-modules of $F$, and let $s$ be a point of $S$. Then there exists an open neighbourhood $U$ of $s$ such that the sequence $(E_n|U)$ is stationary.

\emph{Proof.} Since the local ring $\mathcal O_{S,s}$ is noetherian, there is an integer $n_0$ such that $(E_n)_s=(E_{n_0})_s$ for $n\geq n_0$. Replacing $(E_n)$ by $(E_n/E_{n_0})$ for $n\geq n_0$, one may therefore suppose that $(E_n)_s=0$ for all $n$.

On the other hand, $F_s=M$ is a finite type $\mathcal O_{S,s}$-module. Hence it admits a finite filtration
\[
M=M_0\supset M_1\supset\cdots\supset M_q=0
\]
such that $M_i/M_{i+1}\simeq \mathcal O_{S,s}/\mathfrak p_i$, where $\mathfrak p_i$ is a prime ideal of $\mathcal O_{S,s}$. Since the category of finite type $\mathcal O_{S,s}$-modules is equivalent to the category of germs at $s$ of coherent $\mathcal O_S$-modules, one may suppose, after localizing near $s$, that the $M_i$ (respectively $\mathfrak p_i$) are the fibres at $s$ of coherent sub-$\mathcal O_S$-modules $F_i$ (respectively coherent ideals $\mathcal P_i$ of $\mathcal O_S$), with
\[
F=F_0\supset F_1\supset\cdots\supset F_q=0,
\]
and
\[
F_i/F_{i+1}\simeq \mathcal O_S/\mathcal P_i.
\]
To prove that $(E_n)$ is stationary near $s$ is evidently equivalent to proving that the induced sequences on each term of the filtration $(F_i)$ are stationary. Thus one may suppose that $F=\mathcal O_S$ and that $\mathcal O_{S,s}$ is an integral ring, the $E_n$ forming an increasing sequence of ideals of $\mathcal O_S$ such that $(E_n)_s=0$.

Since $\mathcal O_{S,s}$ is integral, in particular reduced, one may moreover suppose, by Oka's theorem and after localizing at $s$, that $S$ is reduced. We shall show that there is an open neighbourhood $U$ of $s$ such that $E_n|U=0$ for all sufficiently large $n$. For this, let $\widetilde S$ be the normalization of $S$. It is known that $S$ is the analytic spectrum of a finite $\mathcal O_S$-algebra $\widetilde{\mathcal O}_S$ whose fibre at each $t\in S$ is the integral closure of $\mathcal O_{S,t}$. One has
\[
E_n\widetilde{\mathcal O}_S
 =p_*\bigl(E_n\mathcal O_{\widetilde S}\bigr),
\]
where $p:\widetilde S\to S$ is the canonical projection and where $E_n\mathcal O_{\widetilde S}$ is the ideal of $\mathcal O_{\widetilde S}$ canonically defined by $E_n\widetilde{\mathcal O}_S$. It is therefore enough to prove that there is an open neighbourhood $V$ of $s$ in $\widetilde S$ such that $E_n\mathcal O_{\widetilde S}|V=0$ for $n$ sufficiently large. This brings us to the case where $S$ is normal.

Since $(E_n)_s=0$, the closed immersion $i_n:S_n\to S$ defined by $E_n$ is etale at $s$. But $S$ being normal, the set of points where $i_n$ is etale is both open and closed by the lemma below; its image in $S$ is therefore an open and closed part of $S$. Hence the restriction of $E_n$ to the connected component of $s$ in $S$ is zero. Since this connected component is open, $S$ being locally connected, the proof is complete. It remains only to prove the following lemma.

\textbf{Lemma 2.4.8.2.} Let $Y$ be a normal pseudo-analytic space and let $f:X\to Y$ be an unramified morphism. The set of points of $X$ at which $f$ is etale is both open and closed.

\emph{Proof.} It is immediate that the set of points at which $f$ is etale is open, since saying that $f$ is etale at a point means that $f$ is an isomorphism in a neighbourhood of that point. It remains to prove that the set of points at which $f$ is not etale is open. Since an unramified morphism is locally an immersion, one may suppose that $f$ is a closed immersion, defined by a coherent ideal $I$. Let $x$ be a point of $X$ such that $I_x\neq 0$. Let $\mathfrak p$ be a prime ideal of $\mathcal O_{Y,x}$ containing $I_x$. Then
\[
\dim(\mathcal O_{Y,x}/\mathfrak p)>0,
\]
for otherwise one would have $I_x\mathfrak p=0$, since $\mathcal O_{Y,x}$ is integral. Thus $Y$ has positive codimension at $x$, and consequently has positive codimension in a neighbourhood of $x$; the assertion follows.

\section*{3. Complements on quasi-coherent sheaves on schemes}

\textbf{3.0.} In this section we fix a quasi-compact and quasi-separated scheme $S$. The canonical fully faithful inclusion functor
\[
\varphi:\Qcoh(S)\longrightarrow\Mod(S)
\]
extends to a functor, again denoted $\varphi$, from $\D(\Qcoh(S))$ to
\[
\D(S)=\D(\Mod(S)),
\]
which we propose to study.

\textbf{Lemma 3.1.} The category $\Qcoh(S)$ is an abelian category satisfying (AB5), and the quasi-coherent sheaves of finite type form a small family of generators. In particular, $\Qcoh(S)$ has enough injectives.

\emph{Proof.} The assertion about generators follows at once from the fact that every quasi-coherent sheaf on $S$ is the filtered inductive limit of its quasi-coherent subsheaves of finite type. The other assertions are evident.

\textbf{Remark 3.1.1.} The editor does not know whether $\Qcoh(S)$ still has enough injectives if one assumes only that $S$ is quasi-separated.

\textbf{Lemma 3.2.} The functor
\[
\varphi:\Qcoh(S)\longrightarrow\Mod(S)
\]
has a right adjoint
\[
Q:\Mod(S)\longrightarrow\Qcoh(S),
\]
called the \emph{coherer}.

\emph{Proof.} If $S$ is affine, the existence of $Q$ is immediate: one takes
\[
Q=\widetilde{\Gamma_S}.
\]
Let now $f:U\to S$ be an affine open of $S$. It is clear that, for $E\in\Ob\Qcoh(S)$ and $F\in\Ob\Mod(U)$, one has a canonical isomorphism
\begin{equation}\tag{1}
\Hom(E,f_*F)\simeq \Hom\bigl(E,f_*Q(F)\bigr).
\end{equation}
Let $(f_i:U_i\to S)$ be a finite covering of $S$ by affine opens, and, for each $(i,j)$, let $(U_{ijk}\to U_{ij})$ be a finite covering of $U_{ij}=U_i\cap U_j$ by affine opens (the $U_{ij}$ are quasi-compact because $S$ is quasi-separated). Let $F\in\Ob\Mod(S)$. One has an exact sequence
\begin{equation}\tag{2}
0\longrightarrow F\longrightarrow \prod_i f_{i*}(F|U_i)
\longrightarrow \prod_{ijk} f_{ijk*}(F|U_{ijk}).
\end{equation}
If $g:V\to U$ is an inclusion of affine opens of $S$, there is an evident natural arrow
\[
Q(F|U)\longrightarrow g_*Q(F|V).
\]
Consequently (2) gives an arrow
\begin{equation}\tag{3}
\prod_i f_{i*}Q(F|U_i)
\longrightarrow \prod_{ijk} f_{ijk*}Q(F|U_{ijk}),
\end{equation}
whose kernel we denote by $QF$. For $E\in\Ob\Qcoh(S)$, applying $\Hom(E,\cdot)$ to (2) and using (1), we obtain a functorial isomorphism
\[
\Hom(E,F)\simeq \Hom(E,QF).
\]
This proves the existence of $Q$; uniqueness is clear since it is an adjoint.

\textbf{Lemma 3.3.} a) The coherer is an additive functor which sends injective objects to injective objects.

b) Let
\[
RQ:\D^+(S)\longrightarrow\D^+(\Qcoh(S))
\]
be the right derived functor of the coherer. For $E\in\Ob\D(\Qcoh(S))$ and $F\in\Ob\D^+(S)$, there are canonical functorial isomorphisms
\begin{equation}\tag{i}
\RHom_{\Mod}(E,F)\simeq \RHom_{\Qcoh}(E,RQ(F)),
\end{equation}
\begin{equation}\tag{ii}
\Hom_{\Mod}(E,F)\simeq \Hom_{\Qcoh}(E,RQ(F)).
\end{equation}
Here the subscript $\Mod$ (respectively $\Qcoh$) is an abbreviation for $\D(S)$ (respectively $\D(\Qcoh(S))$).

\emph{Proof.} Assertion a) is an immediate consequence of the definition. For b), it is enough to prove the isomorphism (i), since (ii) follows from it by passing to cohomology. We may suppose $F$ has injective components; then $QF$ has injective components by a), and (i) follows directly from the definition of $Q$.

\textbf{Lemma 3.4.} Let $f:X\to Y$ be a morphism of quasi-compact and quasi-separated schemes. For $E\in\Ob\Mod(X)$, there is a canonical functorial isomorphism
\[
Q_Y f_*(E)\simeq f_*Q_X(E).
\]

\emph{Proof.} Since $f$ is quasi-compact and quasi-separated, $f_*$ sends quasi-coherent sheaves to quasi-coherent sheaves. The two functors $Q_Y f_*$ and $f_*Q_X$, from $\Mod(X)$ to $\Qcoh(Y)$, are both right adjoints to
\[
f^*:\Qcoh(Y)\longrightarrow\Mod(X),
\]
hence are canonically isomorphic.

\textbf{Proposition 3.5.} If $S$ is noetherian, or quasi-compact and separated, the functor
\[
\varphi:\D^+(\Qcoh(S))\longrightarrow\D^+(S)
\]
is fully faithful and has as essential image the full subcategory $\D^+(S)_{\qcoh}$ of $\D^+(S)$ formed by complexes with quasi-coherent cohomology.

\emph{Proof.} We may restrict to proving the assertion when $S$ is separated, the case where $S$ is noetherian being well known. It is plainly enough to prove the following two assertions:
\begin{enumerate}[label=(3.5.\arabic*)]
\item for $E\in\Ob\D^+(\Qcoh(S))$, the adjunction arrow
\[
E\longrightarrow RQ(E)
\]
is an isomorphism;
\item for $E\in\Ob\D^+(S)_{\qcoh}$, the adjunction arrow
\[
RQ(E)\longrightarrow E
\]
is an isomorphism.
\end{enumerate}

Proof of (3.5.1). By an immediate devissage, one is reduced to proving that quasi-coherent sheaves are acyclic for the coherer; that is, for $E\in\Ob\Qcoh(S)$ one has
\[
R^iQ(E)=0\quad (i>0),
\]
and the adjunction arrow $E\to QE$ is an isomorphism. If $S$ is affine, this follows from the formula $Q=\widetilde{\Gamma_S}$. In general, cover $S$ by a finite number of affine opens $U_i$. Since $S$ is separated, the inclusions $f_i:U_i\to S$ and also the inclusions
\[
f_{i_0\ldots i_p}:U_{i_0\ldots i_p}=U_{i_0}\cap\cdots\cap U_{i_p}\to S
\]
are affine morphisms. For $E\in\Ob\Qcoh(S)$, consider the resolution of $E$ defined by the covering $(U_i)$:
\begin{equation}\tag{*}
0\longrightarrow E\longrightarrow \prod_i f_{i*}f_i^*E\longrightarrow\cdots\longrightarrow
\prod_{i_0\ldots i_p} f_{i_0\ldots i_p,*}f_{i_0\ldots i_p}^*E\longrightarrow\cdots.
\end{equation}
It gives a spectral sequence
\begin{equation}\tag{**}
E_1^{pq}=\prod_{i_0\ldots i_p} R^qQ\bigl(f_{i_0\ldots i_p,*}f_{i_0\ldots i_p}^*E\bigr)
\Longrightarrow R^{p+q}Q(E).
\end{equation}
Since the $f_{i_0\ldots i_p}$ are affine, one has by (3.4)
\begin{equation}\tag{***}
R^qQ\bigl(f_{i_0\ldots i_p,*}f_{i_0\ldots i_p}^*E\bigr)
\simeq
f_{i_0\ldots i_p,*} R^qQ\bigl(f_{i_0\ldots i_p}^*E\bigr).
\end{equation}
For all $q$; hence in particular
\[
E_1^{pq}=0\quad\text{for }q>0,
\]
by the affine case already treated. From (**) one deduces isomorphisms
\[
R^nQ(E)\simeq H^n(QE')\quad(n\geq0),
\]
where $E'$ is the resolution of $E$ defined by the exact sequence (*). By (***), $QE'\simeq E'$ in positive degrees, and consequently
\[
R^nQ(E)=0\quad(n>0),
\]
and $Q(E)\simeq E$.

Proof of (3.5.2). By the spectral sequence
\[
R^pQ\bigl(H^q(E)\bigr)\Longrightarrow R^{p+q}Q(E),
\]
one is immediately reduced to the case where $E\in\Ob\Qcoh(S)$. The composite
\[
E\longrightarrow RQE\longrightarrow E
\]
of the adjunction arrows is the identity, and the assertion follows from (3.5.1). This completes the proof of (3.5).

\textbf{Remark 3.6.} The conclusion of (3.5) would be false if the hypothesis ``separated or noetherian'' were omitted, as is shown by a useful counterexample of J.-L. Verdier (Appendix I to this exposé).

\textbf{Proposition 3.7.} Under the hypothesis of (3.5), suppose moreover that there is an integer $N$ such that, for every $F\in\Ob\Mod(S)$ and every quasi-compact open $U$ of $S$,
\[
H^i(U,F)=0\quad\text{for }i>N.
\]
Then:
\begin{enumerate}[label=\alph*)]
\item the coherer has finite cohomological dimension;
\item the functor
\[
\varphi:\D(\Qcoh(S))\longrightarrow\D(S)
\]
is fully faithful and has as essential image the full subcategory $\D(S)_{\qcoh}$ of $\D(S)$ formed by complexes with quasi-coherent cohomology.
\end{enumerate}

\emph{Proof of a).} If $S$ is affine, the formula $Q=\widetilde{\Gamma_S}$ shows that $Q$ has cohomological dimension $\leq N$. Passing to the general case, let us prove that there is an integer $a$ such that $R^qQ(E)=0$ for every $E\in\Ob\Mod(S)$ and every $q>a$. The hypothesis implies that every $E\in\Ob\Mod(S)$ admits a finite right resolution $E'$ of length $N$ such that all the $E^i$ are acyclic, i.e. $H^q(U,E^i)=0$ for every open $U$ of $S$ and every $q>0$. Therefore it suffices to prove that there is an integer $b$ such that $R^qQ(E)=0$ for every acyclic $\mathcal O_S$-module $E$ and every $q>b$; then one may take $a=b+N$.

Let $(U_i)$ be a finite covering by $N'$ affine opens. Let $E$ be an acyclic $\mathcal O_S$-module, and let $C^\bullet(E)$ be the resolution of $E$ defined by the covering $(U_i)$. Two cases have to be examined.

First suppose $S$ is separated. The components of $C^\bullet(E)$ are sums of terms of the form $f_*f^*E$, where $f:U\to S$ is the inclusion of an affine open, hence an affine morphism. We have
\[
R(Qf_*)f^*E\simeq RQ\,Rf_*f^*E\simeq RQ\,f_*f^*E,
\]
since $f^*E$ is acyclic. On the other hand, by (3.4),
\[
R(Qf_*)f^*E\simeq R(f_*Q)f^*E\simeq f_*RQ\,f^*E,
\]
since $f$ is affine. Hence
\[
R^nQ(f_*f^*E)\simeq f_*R^nQ(f^*E)\quad(n\geq0),
\]
and the affine case gives
\[
R^nQ(f_*f^*E)=0\quad(n>N).
\]
Since $C^\bullet(E)$ has length $N'$, it follows that
\[
R^nQ(E)=0\quad(n>N+N').
\]

Second suppose $S$ is noetherian. The components of $C^\bullet(E)$ are sums of terms $f_*f^*E$, where $f:U\to S$ is the inclusion of a quasi-compact separated open. As in the first case,
\[
R(Qf_*)f^*E\simeq RQ\,f_*f^*E\simeq R(f_*Q)f^*E.
\]
Since $U$ is noetherian, $\Qcoh(U)$ has enough injective objects in $\Mod(U)$, and the right derived functor of $f_*:\Qcoh(U)\to\Qcoh(S)$ is induced by $Rf_*:\D^+(U)\to\D^+(S)$ when one identifies $\D^+(\Qcoh(U))$ and $\D^+(\Qcoh(S))$ with full subcategories of $\D^+(U)$ and $\D^+(S)$ by (3.5). Therefore
\[
R(f_*Q)f^*E\simeq Rf_*RQ\,f^*E.
\]
This yields a spectral sequence
\[
E_2^{pq}=R^pf_*R^qQ(f^*E)\Longrightarrow R^{p+q}(f_*Q)f^*E.
\]
By the first case, there is an integer $n(U)$, depending only on $U$, such that $Q$ has cohomological dimension $\leq n(U)$ on $\Mod(U)$; hence
\[
E_2^{pq}=0\quad(q>n(U)).
\]
Moreover
\[
E_2^{pq}=0\quad(p>N)
\]
by the cohomological-dimension hypothesis on $S$. Therefore
\[
R^n(f_*Q)f^*E=0\quad(n>N+N''),
\]
where $N''$ is an integer dominating all the $n(U)$, as $U$ runs through the finite set of intersections $U_{i_0}\cap\cdots\cap U_{i_p}$. Taking account of the previous displayed isomorphism, one obtains
\[
R^nQ(f_*f^*E)=0\quad(n>N+N''),
\]
and finally
\[
R^nQ(E)=0\quad(n>N+N'+N''),
\]
which completes the proof of a).

\emph{Proof of b).} Since the coherer has finite cohomological dimension by a), it admits a right derived functor
\[
RQ:\D(S)\longrightarrow\D(\Qcoh(S)),
\]
whose restriction to $\D^+(S)$ is the functor defined above. For $E\in\Ob\D(\Qcoh(S))$, there is a canonical functorial isomorphism
\begin{equation}\tag{1}
E\xrightarrow{\sim} RQ(E),
\end{equation}
since, by (3.5.1), the $E^i$ are acyclic for $Q$ and satisfy $E^i\simeq Q(E^i)$. On the other hand, for $E\in\Ob\D(S)$ there is a canonical functorial homomorphism
\begin{equation}\tag{2}
RQ(E)\longrightarrow E
\end{equation}
coming, when the $E^i$ are $Q$-acyclic, from the ordinary adjunction homomorphism between $\varphi$ and $Q$. It is immediate that (1) and (2) make $RQ:\D(S)\to\D(\Qcoh(S))$ a right adjoint to $\varphi:\D(\Qcoh(S))\to\D(S)$. The fact that (1) is an isomorphism implies that $\varphi$ is fully faithful. Finally, using the spectral sequence
\[
R^pQ\bigl(H^q(E)\bigr)\Longrightarrow R^{p+q}Q(E),
\]
which converges since $Q$ has finite cohomological dimension, one sees that (2) is an isomorphism for $E\in\Ob\D(S)_{\qcoh}$; hence $\varphi$ has essential image $\D(S)_{\qcoh}$. This proves b) and completes the proof of (3.7).

\section*{Appendix I. A counterexample of Verdier}

\subsection*{0. Introduction}
If $S$ is a scheme, write $\Mod(S)$ (respectively $\Qcoh(S)$) for the category of $\mathcal O_S$-modules (respectively quasi-coherent $\mathcal O_S$-modules). The purpose of this appendix is to draw the reader's attention to certain pathological properties of the category $\D(\Qcoh(S))$, and in particular to prove the following facts.

\textbf{0.1.} Let $X$ be an affine scheme with ring $A$, and let $I$ be an injective $A$-module. It is known that, if $A$ is noetherian, $\widetilde I$ is an injective $\mathcal O_X$-module. If $A$ is not noetherian, $\widetilde I$ is not necessarily injective, nor even flasque, even if the underlying space of $X$ is noetherian and of finite Krull dimension. Moreover, the restriction of $\widetilde I$ to an open $U$ of $X$ is not in general an injective object of $\Qcoh(U)$.

\textbf{0.2.} Let $f:X\to Y$ be a morphism of quasi-compact and quasi-separated schemes. The functor
\[
Rf_*:\D^+(X)\longrightarrow\D^+(Y)
\]
does not in general induce a functor on $\D^+(\Qcoh(X))$ with values in $\D^+(\Qcoh(Y))$ (the full subcategories of (II 3.5)), nor is it in general the right derived functor of
\[
f_*:\Qcoh(X)\longrightarrow\Mod(Y).
\]
More precisely, an injective object of $\Qcoh(X)$ is not necessarily acyclic for $f_*:\Mod(X)\to\Mod(Y)$.

\textbf{0.3.} Let $S$ be a quasi-compact and quasi-separated scheme. An object of $\Qcoh(S)$, even an injective one, is not necessarily acyclic for the coherer (II 3.2). The functor
\[
\varphi:\D^b(\Qcoh(S))\longrightarrow\D^b(S)
\]
of (II 3.0) is not in general fully faithful.

\subsection*{1. The counterexamples}

First recall the following lemma (SGA 2 II 9).

\textbf{Lemma 1.1.} Let $A$ be a ring, $X=\Spec(A)$, let $\mathfrak f=(f_\alpha)$ be a finite system of elements of $A$, and let $Y$ be the closed subset defined by $\mathfrak f(x)=0$. For $i>0$, the following conditions are equivalent:
\begin{enumerate}[label=(\roman*)]
\item for every injective $A$-module $I$, one has $H_Y^i(X,\widetilde I)=0$;
\item the pro-object
\[
n\longmapsto H_i(\mathfrak f^n,A)
\]
is zero, that is, for every $n$ there exists $n'\geq n$ such that
\[
H_i(\mathfrak f^{n'},A)\longrightarrow H_i(\mathfrak f^n,A)
\]
is zero.
\end{enumerate}

\textbf{1.2.} We shall use (1.1) to prove (0.1). Let $k$ be a field,
\[
B=k[[x,y]],\qquad \mathfrak f=(x,y),\qquad
M=\coprod_{n\geq0} B/(\mathfrak f^n),
\]
and let $A$ be the $B$-algebra $D_B(M)$. One has trivially
\[
H^A_2(\mathfrak f^n,A)\simeq H^B_2(\mathfrak f^n,A).
\]
On the other hand
\[
H^B_2(\mathfrak f^n,A)\simeq H^B_2(\mathfrak f^n,B)\oplus H^B_2(\mathfrak f^n,M),
\]
but $H^B_2(\mathfrak f^n,B)=0$, since $\mathfrak f$ is a regular sequence; hence
\[
H^B_2(\mathfrak f^n,A)\simeq H^B_2(\mathfrak f^n,M).
\]
Finally,
\[
H^B_2(\mathfrak f^n,M)\simeq H^0_B(\mathfrak f^n,M),
\]
where $H^0_B(\mathfrak f^n,M)$ identifies canonically with the submodule $(\mathfrak f^n)^M$ of $M$ annihilated by $\mathfrak f^n$. Thus
\begin{equation}\tag{*}
H^A_2(\mathfrak f^n,A)\simeq (\mathfrak f^n)^M.
\end{equation}
Since the transition morphism
\[
H^A_2(\mathfrak f^{n'},A)\longrightarrow H^A_2(\mathfrak f^n,A)
\]
is induced by multiplication by $\mathfrak f^{n'-n}$, one verifies immediately from (*) that the pro-object $n\mapsto H^A_2(\mathfrak f^n,A)$ is not zero. It follows from (1.1) that there exists an injective $A$-module $I$ such that $H_Y^2(X,\widetilde I)\neq0$, hence such that $\widetilde I$ is not flasque. Notice also that the projection
\[
\Spec(A)\longrightarrow\Spec(B)
\]
is a homeomorphism, so that the underlying space of $X$ is noetherian of dimension $2$. Let $U$ be the open complement of $Y$ in $X$. One has
\[
H_Y^2(X,\widetilde I)\simeq H^1(U,\widetilde I)\neq0.
\]
Since $U$ is separated, the functor $\D^+(\Qcoh(U))\to\D^+(U)$ is fully faithful by II 3.5; hence $\widetilde I|U$ is not injective in $\Qcoh(U)$. This proves (0.1).

Let $Z$ be the quasi-compact and quasi-separated scheme obtained by gluing two copies of $X$ along $U$, and let $j_1,j_2$ be the two canonical open immersions of $X$ into $Z$. Thus one has a Cartesian square
\[
\begin{tikzcd}
U \arrow[r,"i"] \arrow[d,"i"'] & X \arrow[d,"j_1"]\\
X \arrow[r,"j_2"'] & Z .
\end{tikzcd}
\]
Then
\[
j_2^*R^1j_{1*}(\widetilde I)\simeq R^1i_*\bigl(i^*\widetilde I\bigr),
\]
whence
\[
H^0\bigl(X,R^1j_{1*}(\widetilde I)\bigr)\simeq H^1(U,\widetilde I)\neq0.
\]
Thus $R^1j_{1*}(\widetilde I)\neq0$, which proves (0.2).

On the other hand, by (II 3.4 and 3.5),
\[
RQ\,Rj_{1*}(\widetilde I)\simeq j_{1*}\widetilde I.
\]
It follows from the spectral sequence and from the exact sequence of low-degree terms that
\[
R^1j_{1*}(\widetilde I)\simeq R^2Q(j_{1*}\widetilde I).
\]
Since the sheaf $j_{1*}\widetilde I$ is injective in $\Qcoh(Z)$, the first part of (0.3) is proved.

Finally, let $E\in\Ob\Qcoh(Z)$. There is the duality isomorphism
\[
\RHom_{\Mod}\bigl(E,Rj_{1*}(\widetilde I)\bigr)
\simeq
\RHom_{\Mod}\bigl(j_1^*E,\widetilde I\bigr),
\]
and, since $X$ is affine, II 3.5 gives
\[
\RHom_{\Mod}(j_1^*E,\widetilde I)
\simeq
\RHom_{\Qcoh}(j_1^*E,I)
\simeq
\Hom(j_1^*E,I).
\]
Therefore
\[
\RHom_{\Mod}\bigl(E,Rj_{1*}(\widetilde I)\bigr)
\simeq
\Hom(j_1^*E,I),
\]
and the resulting spectral sequence gives the low-degree isomorphism
\[
\Ext^2_{\Mod}\bigl(E,j_{1*}\widetilde I\bigr)
\simeq
\Hom\bigl(E,R^1j_{1*}\widetilde I\bigr).
\]
Consequently the morphism $\Id_E$ of $E=R^1j_{1*}\widetilde I$ provides a nonzero element of $\Ext^2_{\Mod}(E,j_{1*}\widetilde I)$; hence the functor
\[
\D^b(\Qcoh(Z))\longrightarrow\D^b(Z)
\]
is not fully faithful. This completes the proof of (0.3).

\end{document}
